\begin{sol}[Roca esférica empujada por una topadora cuesta arriba]

\subsubsection*{Marcos y cinemática}

Se usa un marco inercial fijo a la ladera con $\vec{E}_x$ cuesta arriba (a lo largo de la
pendiente), $\vec{E}_y$ normal a la superficie y $\vec{E}_z=\vec{E}_x\times\vec{E}_y$. El
centro de la esfera está en $\vec{r}_O=x\vec{E}_x$, de modo que $\vec{a}_O=\ddot{x}\vec{E}_x$.
La condición de rodadura sin deslizar en el punto de contacto $Q$ con la ladera da
$$\dot{x}=r\omega \;\Longrightarrow\; \omega=\frac{\dot{x}}{r},\qquad \dot\omega=\frac{\ddot{x}}{r}.$$

\subsubsection*{Fuerzas}

El punto de la esfera en contacto con la hoja ($P$) desliza respecto a ésta, por lo que la
fricción dinámica de la hoja es $\vec{F}_f=\mu P\,\vec{E}_y$ (magnitud $\mu P$). Además actúan
la fuerza $\vec{P}=P\vec{E}_x$, el peso $\vec{W}=-mg\sin\beta\,\vec{E}_x + mg\cos\beta\,\vec{E}_y$,
y la normal y la tracción del suelo en $Q$ (que no generan torque respecto a $Q$).

\subsubsection*{Balance de momento angular respecto al contacto $Q$}

Para evitar las reacciones desconocidas del suelo, se toma momento angular respecto al punto
de contacto $Q$. Con $\bar{I}=\tfrac25 mr^2$:
$$\vec{H}_Q = \bar{I}\,\omega\,\vec{E}_z + (\vec{r}_O-\vec{r}_Q)\times m\vec{v}_O
= \frac{2}{5}mr\dot{x}\,\vec{E}_z + mr\dot{x}\,\vec{E}_z = \frac{7}{5}mr\dot{x}\,\vec{E}_z,$$
de donde $\dot{\vec{H}}_Q = \tfrac75 mr\ddot{x}\,\vec{E}_z$.

El torque respecto a $Q$, con $\vec{r}_P-\vec{r}_Q=-r\vec{E}_x-r\vec{E}_y$ y
$\vec{r}_O-\vec{r}_Q=-r\vec{E}_y$, resulta
$$\vec{M}_Q = (\vec{r}_P-\vec{r}_Q)\times(\vec{P}+\vec{F}_f) + (\vec{r}_O-\vec{r}_Q)\times\vec{W}
= \big[rP(1-\mu) - mgr\sin\beta\big]\vec{E}_z.$$

Igualando $\vec{M}_Q=\dot{\vec{H}}_Q$ y dividiendo por $r$:
$$rP(1-\mu) - mgr\sin\beta = \frac{7}{5}mr\ddot{x}
\;\Longrightarrow\;
\boxed{\;\frac{7}{5}m\ddot{x} + mg\sin\beta = P(1-\mu)\;}$$

Equivalentemente, $\ddot{x} = \dfrac{5}{7m}\big[P(1-\mu) - mg\sin\beta\big]$.

\end{sol}
